 \part{Model checking with type trees}
 \label{part:mc}

 In this part we prove:
{\renewcommand{\mcdepth}{d}\mcthm*}

\section{Preliminaries}\label{sec:mc-prelims}

In \Cref{part:mc},
unless stated otherwise, $G$ denotes a partial graph.
We slightly extend the notions of partial graphs and of construction sequences.

% lax begin Lax185574.PartialGraphs
\paragraph{Partial graphs.}
To track parts over time, we give them names.
Fix a countably infinite set of unary symbols, called \emph{colors}.
We regard a partial graph $G$ as a relational structure with the two symmetric binary relations $E$ and $N$ and a unary relation $Y$ for every color $Y$.
The vertices of color $Y$ form its color class $Y^G$, and the nonempty color classes are the parts of $G$.
We write $|G|\coloneqq|V(G)|$ and $\size G\coloneqq|V(G)|+|R(G)|$ for the size of the representation of a partial graph $G$.
% lax end
% lax begin Lax185574.DistanceMatrix
The \emph{distance matrix} of $G$ is the $|G|\times|G|$ matrix $D_G$ with entries $D_G(u,v)\coloneqq\dist^G(u,v)\in\N\cup\set\infty$.
Passing it as an additional input to the algorithms below saves time, as distances need not be recomputed repeatedly.
% lax end

% lax begin Lax185574.ConstructionSequences
\paragraph{Construction sequences.}
For colors $P,Q$ and a partial graph $G$, we define three partial graphs
$\merge PQ(G)$ (for $P\neq Q$), $\resolve+PQ(G)$, and $\resolve-PQ(G)$, as the partial graph $H$ obtained from $G$ by leaving all relations of $G$ unchanged, apart from redefining the following:
\begin{align*}
\merge PQ(G)\colon\quad&
P^{H}\coloneqq P^G\cup Q^G,
\qquad
Q^{H}\coloneqq\emptyset
&&\text{for $P\ne Q$},\\
\resolve+PQ(G)\colon\quad&
E(H)\coloneqq E(G)\cup\bigl(P^GQ^G\setminus N(G)\bigr),\\
\resolve-PQ(G)\colon\quad&
N(H)\coloneqq N(G)\cup\bigl(P^GQ^G\setminus E(G)\bigr).
\end{align*}
After either resolve, every pair in $P^GQ^G$ is resolved.
In this representation, every step of a construction sequence $G_1,\ldots,G_\lambda$ (\Cref{sec:tech-overview}) has the form $G_{t+1}=\merge PQ(G_t)$, $\resolve+PQ(G_t)$, or $\resolve-PQ(G_t)$ for some colors $P,Q$.
% lax end

\section{Distance logic}\label{sec:dist-rank}

Following \cite{mw-arxiv,mw-stoc} (and earlier work by \cite{gks}),
we define a logic called \distFO, extending first-order logic by distance atoms, and a suitable notion of rank. We then define compact representations of types tailored to this logic.

\subsection{Distance logic and distance rank}

% lax begin Lax185574.DistanceLogic
\paragraph{First-order distance logic.}
The logic \distFO extends first-order logic by additionally allowing the following two types of atoms:
\begin{quote}
     For each radius \(r \in \N\), \distFO introduces
        a \emph{binary distance atom} \(\dist(x,y) \le r\) expressing that the distance between \(x\) and \(y\) is at most \(r\).

     For each radius \(r \in \N\) and color \(Y\), \distFO introduces
        a \emph{unary distance atom} \({\dist(x,Y) < r}\) expressing that the distance from \(x\) to \(Y\) is smaller than \(r\),
        equivalent to \({\exists y~Y(y) \land \dist(x,y) < r}\).
\end{quote}
In both cases, we call \(r\) the \emph{radius} of the distance atom.
Distance atoms are expressible in first-order logic, but only with quantifiers.
Treating them as atoms lets us measure distances without increasing the quantifier rank, as in~\cite{gks}.

\paragraph{Horizon functions.}
We fix sufficiently fast-growing horizon functions as in \cite{locality-note}:
\begin{equation}\label{eq:horizon-def}
\rhom(k,q)\coloneqq 9^{(k+q+1)q}
\qquad\text{and}\qquad
\rhop(k,q)\coloneqq 9^{(k+q)(q+1)}.
\end{equation}
They satisfy the two identities
\begin{equation}\label{eq:horizon}
\rhop(k,q)=9^k\rhom(k,q)\quad\text{for all $k,q\in\N$},
\qquad
\rhom(k,q)=9^q\rhop(k+1,q-1)\quad\text{for $q\ge1$}.
\end{equation}
In particular, as $9^q\ge2$,
\begin{equation}\label{eq:horizon-cons}
\rhop(k,q)\ge\rhom(k,q)\ge2\rhop(k+1,q-1)\quad\text{for all $k\in\N$ and $q\ge1$}.
\end{equation}

\paragraph{Distance rank.}
Let $k,q\in\N$.
A \distFO formula has \emph{distance rank \((k,q)\)} if it is a Boolean combination of
\begin{itemize}
\item first-order atoms: relational symbols in the signature or equality, applied to variables,
\item binary distance atoms with radius $r\le\rhom(k,q)$,
\item unary distance atoms with radius $r\le\rhop(k,q)$,
\item
formulas
\[
\exists y~\Bigl(\dist(x_i,y)\le \rhop(k+1,q-1)\land\phi(\tup x y)\Bigr),
\]
where \(q \ge 1\), $x_i$ is a variable in $\tup x$, and $\phi(\tup x y)$ has distance rank $(k+1,q-1)$.
\end{itemize}

A formula with distance rank \((k,q)\) has quantifier rank at most \(q\).
Observe by \eqref{eq:horizon-cons} that
the distance atoms have progressively shorter radii as one traverses the quantifiers
inwards.

Formulas of distance rank $(k,q)$ are called \emph{local} formulas of distance rank $(k,q)$ in
\cite{mw-arxiv,locality-note}, as opposed to the more general formulas which also allow the unrestricted quantification $\exists y\, \psi(\tup x y)$. Since we only consider local formulas in this paper, we omit this adjective.

% lax end

\subsection{Type trees}
\label{sec:typetrees}
 In the context of usual first-order logic, the \emph{type} of quantifier rank $q$ of a tuple $\tup a$ in a structure $A$ is defined as the set of all formulas of quantifier rank $q$ that are satisfied by $\tup a$ in $A$ (see \Cref{sec:prelims}). It is well known~\cite{Scott1979}
  that for structures over a finite relational signature, such types can be identified with \emph{trees}, also called \emph{Hintikka types},
defined recursively as follows:
$$\tree_q^A(\tup a)=\begin{cases}
\atp^A(\tup a)&\text{if $q=0$,}\\
\setof{\tree_{q-1}^A(\tup a b)}{b\in V(A)}&\text{if $q>0$}.
\end{cases}$$
In the setting of \distFO formulas of distance rank $(k,q)$, the (local) \emph{type of distance rank $(k,q)$} of a tuple $\tup a$
is defined  \cite{mw-arxiv} as the set of (local) \distFO formulas of distance rank $(k,q)$ satisfied by $\tup a$.
In this section, we adapt the construction of trees to this setting. This is useful in the algorithmic context, as trees are much smaller objects than types.

% lax begin Lax185574.TypeTrees
\paragraph{Type trees.}
In the two definitions below, let $G$ be a partial graph, $k,q\in\N$, and $\tup a\in V(G)^{\le k}$, and put $\ell\coloneqq|\tup a|$.
The \emph{atomic $(k,q)$-type}, denoted $\atp_{k,q}(\tup a)$, is the set of all atomic formulas $\phi$
of distance rank \((k,q)\)
with free variables among $x_1,\ldots,x_\ell$, such that $G\models\phi(\tup a)$. Note that this set is finite, as an atom mentioning a color that does not occur in $G$ is not satisfiable in~$G$.

Define the \emph{$(k,q)$-type tree} of $\tup a$ as
\[
\tree_{k,q}\bigl(\tup a\bigr)
\coloneqq
\bigl(\atp_{k,q}(\tup a),\mathcal E_{k,q}(\tup a)\bigr)
\]
where $\mathcal E_{k,q}(\tup a)=\emptyset$ if $q=0$, and otherwise
\[
\mathcal E_{k,q}\bigl(\tup a\bigr)
\coloneqq
\setof{\bigl(i,\tree_{k+1,q-1}(\tup a b)\bigr)}{i\in[\ell],\ b\in V(G),\ \dist(a_i,b)\le\rhop(k+1,q-1)}.
\]

Note that type trees are ``local'' in the sense that new vertices are always quantified close to previous ones (see \Cref{lem:locality}).
The members of $\mathcal E_{k,q}(\tup a)$ are the \emph{extensions} of $\tree_{k,q}(\tup a)$. The position $i$ of an extension $(i,\mu)$ is called the \emph{guard}, and we say that $i$ \emph{admits} $b$ if $\dist(a_i,b)\le\rhop(k+1,q-1)$.
We view $\tree_{k,q}(\tup a)$ as a labeled, rooted tree
 of depth $q$,
with the root labeled $\atp_{k,q}(\tup a)$ and the root of each type tree $\mu$ of an extension $(i,\mu)\in \mathcal E_{k,q}(\tup a)$ attached as a child of the root via an edge labeled $i$.
We call $\ell$ the \emph{arity} of the type tree, and $(k,q)$ its \emph{rank}.

\paragraph{Representation.}
Type trees are represented in the standard RAM model in a canonical (say, sorted,
duplicate-free) form, so that equality of two type trees $\tau_1,\tau_2$ can be
tested in time $O(\size{\tau_1}+\size{\tau_2})$.
By $\size\tau$ we denote the size of this representation (which includes $k,q$, and $\ell$).
% lax end

% lax begin Lax185574.CanonicalForm
% lax begin Lax185574Proofs.CanonicalForm.exists_program_canonical
\begin{observation}[canonical form]\label{obs:canonical}
A type tree given by a representation whose sets need not be sorted or duplicate-free can be brought into canonical form in polynomial time.
\end{observation}
% lax end
% lax end

The next lemma states that the type trees indeed represent the types. It is the reason for computing type trees in the first place.

% lax begin Lax185574.Evaluation
\begin{lemma}[evaluation]\label{lem:eval}
Let $\phi(\tup x)$ be a \distFO formula of distance rank $(k,q)$, \(G\) be a partial graph and $\tup a\in V(G)^{\ell}$ with $\ell=|\tup x|\le k$.
Then whether $G\models\phi(\tup a)$ holds is decidable in polynomial time, given $\tree^G_{k,q}(\tup a)$ and $\phi$.
\end{lemma}
% lax end

\textbf{Notation.} For a $(k,q)$-type tree $\tau=\tree^G_{k,q}(\tup a)$
and a \distFO formula $\phi(\tup x)$ of distance rank $(k,q)$,
we write $\tau \models\phi$, if $G\models\phi(\tup a)$.
The lemma implies that this does not depend on $G$ and $\tup a$, only on the type tree $\tau$.

% lax begin Lax185574Proofs.Evaluation.exists_program_eval
\begin{proof}[Proof of \Cref{lem:eval}]
By induction on $q$.
Write $\tau\coloneqq\tree_{k,q}(\tup a)$.
By definition of distance rank, $\phi$ is a Boolean combination of atoms and of existential components, and we evaluate these components separately.
An atom holds of $\tup a$ if and only if it occurs in the root label $\atp_{k,q}(\tup a)$ of $\tau$.
Now let $q\ge1$ and put $H\coloneqq\rhop(k+1,q-1)$.
Every existential component $\chi$ of $\phi$ can be written as $$\exists y~(\dist(x_i,y)\le H\land\psi(\tup x y))$$ with $i\in[\ell]$ and $\psi$ of distance rank $(k+1,q-1)$.
By the inductive assumption, whether $G\models\psi(\tup c)$ holds is determined by $\tree_{k+1,q-1}(\tup c)$ and $\psi$, and by the notation above, this holds if $\tree_{k+1,q-1}(\tup c)\models \psi$.
Unfolding the semantics, then applying the inductive assumption, and finally the definition of $\mathcal E_{k,q}(\tup a)$, we obtain
\begin{align*}
G\models\chi(\tup a)
&\iff\exists b\in V(G)\colon\ \dist(a_i,b)\le H\ \text{and}\ G\models\psi(\tup ab)\\
&\iff\exists b\in V(G)\colon\ \dist(a_i,b)\le H\ \text{and}\ \tree_{k+1,q-1}(\tup ab)\models\psi\\
&\iff\exists\mu\colon\ (i,\mu)\in\mathcal E_{k,q}(\tup a)\ \text{and}\ \mu\models\psi.
\end{align*}
So $\chi$ is decided by evaluating $\psi$, via the inductive assumption, on the extensions of $\tau$ with guard~$i$.

This proof yields a polynomial-time algorithm.
\end{proof}
% lax end

% lax begin Lax185574.TreeEquivalence
\begin{corollary}\label{cor:tp-tree}
   Fix $\ell\le k$ and $q$,  a partial graph $G$, and  $\tup a,\tup b\in V(G)^\ell$. Then $\tree^G_{k,q}(\tup a)=\tree^G_{k,q}(\tup b)$ if and only if $G\models \phi(\tup a)\Leftrightarrow \phi(\tup b)$ for every \distFO formula $\phi(\tup x)$  of distance rank $(k,q)$ with $|\tup x|=\ell$.
\end{corollary}
% lax end
% lax begin Lax185574Proofs.TreeEquivalence.tree_eq_iff
\begin{proof}
   The rightwards implication follows from \Cref{lem:eval}. The converse is proved by an easy induction on $q$.
\end{proof}
% lax end

% lax begin Lax185574.TreeLocality
\begin{lemma}[locality]\label{lem:locality}
Let $k\ge1$, $q\in\N$, let $G$ be a partial graph, and let $a\in V(G)$.
Every node at depth $j$ of $\tree_{k,q}(a)$ is $\tree_{k+j,q-j}(\tup v)$ for some tuple $\tup v$ from $\Ball_{\rhop(k,q)}(a)$, and every color mentioned in the root label intersects this ball.
\end{lemma}
% lax end
% lax begin Lax185574Proofs.TreeLocality.ball_of_guarded
\begin{proof}
By the definition of extensions and induction on $j$, a node at depth $j$ is $\tree_{k+j,q-j}(\tup v)$ for a tuple $\tup v=(v_0,\ldots,v_j)$ with $v_0=a$ in which every $v_i$ with $i\ge1$ is within distance $\rhop(k+i,q-i)$ of an earlier entry.
By induction on $i$ and the triangle inequality, $\dist(a,v_i)\le\sum_{i'=1}^{i}\rhop(k+i',q-i')$.
The root label mentions colors only in atoms $Y(x_i)$ and $\dist(x_i,Y)<r$ with $r\le\rhop(k+j,q-j)$, and such an atom holds of $\tup v$ only if a vertex of $Y$ is within distance $r$ of $v_i$.
Hence every entry of $\tup v$, and some vertex of every color in an atom in the root label, is within distance
\[
\sum_{i=1}^{j}\rhop(k+i,q-i)+\rhop(k+j,q-j)\le\rhop(k,q)
\]
of $a$, where the inequality uses $\rhop(k+i,q-i)\le2^{-i}\rhop(k,q)$, which follows by applying \eqref{eq:horizon-cons} $i$ times.
\end{proof}
% lax end

\subsection{Operations on type trees}

% lax begin Lax185574.Restriction
\begin{lemma}[ordered coordinate restriction]\label{lem:restrict}
There is an operation $\restriction$ on type trees
such that for all \(k,q\in\N\), partial graphs \(G\),
tuples $\tup w\in V(G)^{\ell}$ for $\ell\le k$ and ordered lists $J$ of distinct positions in $[\ell]$,
\[
\tree_{k,q}(\tup w)\restriction J
=
\tree_{k,q}(\tup w\restriction J).
\]
Given $\tau$ and $J$ one can compute $\tau\restriction J$ in time polynomial in $\size\tau$.
\end{lemma}
% lax end
% lax begin Lax185574Proofs.Restriction.exists_program_restrict
\begin{proof}
By induction on $q$.
Let $J=(j_1,\ldots,j_p)$, and $H\coloneqq\rhop(k+1,q-1)$.
Renaming variables gives
\[
\atp_{k,q}(\tup w\restriction J)=\setof{\psi(x_1,\ldots,x_p)}{\psi(x_{j_1},\ldots,x_{j_p})\in\atp_{k,q}(\tup w)},
\]
which settles the root label.
Let $q\ge1$, and put $J^+\coloneqq(j_1,\ldots,j_p,\ell+1)$, so that $(\tup w\restriction J)\,b=(\tup w b)\restriction J^+$ for every vertex $b$.
Unfolding the definition of $\mathcal E_{k,q}(\tup w\restriction J)$, then applying the inductive assumption, and finally the definition of $\mathcal E_{k,q}(\tup w)$, we obtain
\begin{align*}
\mathcal E_{k,q}(\tup w\restriction J)
&=\setof{\bigl(i,\ \tree_{k+1,q-1}((\tup w b)\restriction J^+)\bigr)}{i\in[p],\ b\in V(G),\ \dist(w_{j_i},b)\le H}\\
&=\setof{\bigl(i,\ \tree_{k+1,q-1}(\tup w b)\restriction J^+\bigr)}{i\in[p],\ b\in V(G),\ \dist(w_{j_i},b)\le H}\\
&=\setof{(i,\ \mu\restriction J^+)}{i\in[p],\ (j_i,\mu)\in\mathcal E_{k,q}(\tup w)}.
\end{align*}
So the guard $j_i$ is renumbered to $i$, as the variable $x_{j_i}$ is renamed to $x_i$ in the root.
The right-hand side depends only on $\tree_{k,q}(\tup w)$ and $J$ and we define $\tau\restriction J$ accordingly.
It is computed by one pass over $\tau$, followed by restoring canonical form. As no subtree grows, this takes polynomial time by \Cref{obs:canonical}.
\end{proof}
% lax end

% lax begin Lax185574.RankShift
\begin{lemma}[rank shift]\label{lem:shift}
There is a polynomial-time algorithm which maps
$\tree_{k,q}(\tup a)$ to $\tree_{k+1,q-1}(\tup a)$ for all \(k\in\N\), \(q\ge1\),
partial graphs \(G\) and nonempty $\tup a\in V(G)^{\le k}$.
\end{lemma}
% lax end
% lax begin Lax185574Proofs.RankShift.exists_program_shift
\begin{proof}
Since $\dist(a_1,a_1)=0$, the guard $1$ admits $a_1$, so $\tree_{k,q}(\tup a)$ has an extension $(1,\mu)$.
Every such $\mu$ is $\tree_{k+1,q-1}(\tup ab)$ for some vertex $b$, and applying the algorithm of \Cref{lem:restrict} yields
$\mu\restriction(1,\ldots,|\tup a|)=\tree_{k+1,q-1}(\tup a)$.
\end{proof}
% lax end

The following is an adaptation of \cite[Lemma~3.1]{locality-note}.
% lax begin Lax185574.FarComposition
\begin{lemma}[far composition]\label{lem:compose}
There is an operation $\ast$ on type trees
such that for all \(k,q\in\N\), partial graphs \(G\) and
nonempty tuples $\tup a_1,\tup a_2$ with $|\tup a_1\tup a_2|\le k$ and $\dist(\tup a_1,\tup a_2)>\rhom(k,q)$,
\[
\tree_{k,q}(\tup a_1)\ast\tree_{k,q}(\tup a_2)
=
\tree_{k,q}(\tup a_1\tup a_2).
\]
Given $\tau_1,\tau_2$ one can compute $\tau_1\ast\tau_2$ in time $\poly_{k,q}(\size{\tau_1}+\size{\tau_2})$.
\end{lemma}
% lax end
% lax begin Lax185574Proofs.FarComposition.exists_program_compose
\begin{proof}
By induction on $q$.
Write $\ell_i\coloneqq|\tup a_i|$ and $\tau_i\coloneqq\tree_{k,q}(\tup a_i)$. Note that \(\ell_i\) can be deduced from~$\tau_i$.

Consider an atom of distance rank $(k,q)$ that mentions variables of both $\tup a_1$ and $\tup a_2$.
If the atom was true, we would have $\dist(\tup a_1,\tup a_2)\le1$ for an equality, edge, or non-edge atom, and $\dist(\tup a_1,\tup a_2)\le\rhom(k,q)$ for a binary distance atom, each contradicting $\dist(\tup a_1,\tup a_2)>\rhom(k,q)$.
Every other atom mentions only variables of $\tup a_1$ or only variables of $\tup a_2$, and whether it holds is recorded in $\atp_{k,q}(\tup a_1)$ or $\atp_{k,q}(\tup a_2)$.
Hence $\atp_{k,q}(\tup a_1\tup a_2)$ is determined by $\atp_{k,q}(\tup a_1)$ and $\atp_{k,q}(\tup a_2)$.
This settles the root label.

Let now \(q \ge 1\).
Put $H\coloneqq\rhop(k+1,q-1)$, and let $\ast$ also denote the operation given by the inductive assumption for rank $(k+1,q-1)$.
Write $\tup a_1\tup a_2=(c_1,\ldots,c_{\ell_1+\ell_2})$, and split $\mathcal E_{k,q}(\tup a_1\tup a_2)$ by guards into $\mathcal E_1\cup\mathcal E_2$, where
\begin{align*}
\mathcal E_1&\coloneqq\setof{\bigl(i,\ \tree_{k+1,q-1}(\tup a_1\tup a_2 b)\bigr)}{i\in[\ell_1],\ b\in V(G),\ \dist(c_i,b)\le H},\\
\mathcal E_2&\coloneqq\setof{\bigl(i,\ \tree_{k+1,q-1}(\tup a_1\tup a_2 b)\bigr)}{\ell_1<i\le\ell_1+\ell_2,\ b\in V(G),\ \dist(c_i,b)\le H}.
\end{align*}

Consider $\mathcal E_1$, and let $i,b$ satisfy its conditions.
Then $\dist(\tup a_1,b)\le H$, and the triangle inequality, the hypothesis, \eqref{eq:horizon-cons} and \eqref{eq:horizon} give
\[
\dist(\tup a_1b,\tup a_2)
\;\ge\;\dist(\tup a_1,\tup a_2)-\dist(\tup a_1,b)
\;>\;\rhom(k,q)-H
\;\ge\;H
\;\ge\;\rhom(k+1,q-1).
\]
Hence, we may apply the inductive assumption for rank $(k+1,q-1)$ to the tuples $\tup a_1b$ and $\tup a_2$.
Moreover, let $J$ be the list of positions with $\tup a_1\tup a_2 b=(\tup a_1 b\tup a_2)\restriction J$.
Applying in this order the property of \(J\), \Cref{lem:restrict} and the inductive assumption, and finally the definition of $\mathcal E_{k,q}(\tup a_1)$, we obtain
\begin{align*}
\mathcal E_1
&=\setof{\bigl(i,\ \tree_{k+1,q-1}(\tup a_1 b\tup a_2)\restriction J\bigr)}{i\in[\ell_1],\ b\in V(G),\ \dist(c_i,b)\le H}\\
&=\setof{\bigl(i,\ \bigl(\tree_{k+1,q-1}(\tup a_1b)\ast\tree_{k+1,q-1}(\tup a_2)\bigr)\restriction J\bigr)}{i\in[\ell_1],\ b\in V(G),\ \dist(c_i,b)\le H}\\
&=\setof{\bigl(i,\ \bigl(\mu\ast\tree_{k+1,q-1}(\tup a_2)\bigr)\restriction J\bigr)}{(i,\mu)\in\mathcal E_{k,q}(\tup a_1)},
\end{align*}
where $\tree_{k+1,q-1}(\tup a_2)$ is obtained from $\tau_2$ by \Cref{lem:shift}.
Symmetrically, with the guard $i$ of $\tup a_2$ renumbered to $\ell_1+i$ as the variables of its root are shifted by $\ell_1$, and without reordering since $\tup a_1\tup a_2 b$ already ends with $b$,
\[
\mathcal E_2
=\setof{\bigl(\ell_1+i,\ \tree_{k+1,q-1}(\tup a_1)\ast\mu\bigr)}{(i,\mu)\in\mathcal E_{k,q}(\tup a_2)}.
\]
The right-hand sides depend only on $\tau_1$ and $\tau_2$ and we define $\tau_1\ast\tau_2$ accordingly.
By \Cref{lem:restrict,lem:shift} and the inductive assumption, the work per input extension is $\poly_{k,q}(\size {\tau_1}+\size{\tau_2})$, and the recursion has fixed depth $q$. Restoring canonical form by \Cref{obs:canonical} stays within this bound, which proves the stated running time.
\end{proof}
% lax end

\section{Scatter types and far trees}\label{sec:localglobal}

The type tree of a tuple $\tup a$ records extensions only by vertices near $\tup a$.
A resolve, however, can bring vertices far from $\tup a$ within guard range, so lifting type trees through a resolve (\Cref{sec:merge-resolve}) also requires the type trees of such far away vertices.
For a vertex $b$ far from $\tup a$, the type tree of $\tup ab$ is obtained from those of $\tup a$ and $b$ by far composition (\Cref{lem:compose}).
As for quantifiers over far witnesses in the proof of Gaifman's locality theorem, it therefore suffices to know which type trees are realized far from $\tup a$.
We collect them in the \emph{far tree} of $\tup a$.

The far tree is not determined by the type tree of $\tup a$ alone. It also depends on global information about $G$.
As in Gaifman's proof, this global information is captured by an analogue of the basic local sentences, which we call the \emph{scatter type} of $G$.
In this section we show that the far tree of $\tup a$ is computable from the type tree of $\tup a$ and the scatter type of $G$.

\Cref{sec:stp} defines scatter types, and \Cref{sec:ftree} uses them to compute the far tree.
The model checking algorithm maintains the type trees of all vertices, and \Cref{lem:compute-stp} shows that the scatter type can be computed from these.

\subsection{Scatter types}\label{sec:stp}

% lax begin Lax185574.ScatterTypes
Let $G$ be a partial graph.
A set $S$ of vertices of $G$ is \emph{$r$-scattered} if all vertices in $S$ have pairwise distance larger than $r$ in $G$.
Fix, once and for all, a linear order on $V(G)$. The algorithms below use it without further mention.
For a set $X\subseteq V(G)$ and $r\in\N$, let $S_r(X)$ denote the inclusion-wise maximal $r$-scattered subset of $X$ constructed by the greedy process that repeatedly adds the least available vertex.
% lax begin Lax185574.TypeTrees
If $\nu$ is a type tree of distance rank $(k,q)$ and arity one, write
\[
[\nu]^G
\coloneqq
\setof{b\in V(G)}{\tree_{k,q}(b)=\nu}
\]
for the set of its realizations.
% lax end

\begin{definition}[scatter type]\label{def:stp}
Let $G$ be a partial graph and $k,q\in\N$. The \emph{$(k,q)$-scatter type} $\stp_{k,q}^G$ is the function  with domain
\begin{align*}
\bigl\{(r,\tree_{k+i,q-i}(b))\ \big|\
i\in[q],\ b\in V(G),\ r\in[9^{k+i}\rhom(k+i,q-i)]\bigr\}
\end{align*}
defined by
\[
\stp^G_{k,q}(r,\nu)
\coloneqq
\min\bigl(|S_r([\nu]^G)|,\ k+q\bigr).
\]
For $q=0$ the domain is empty and $\stp_{k,0}$ is the empty function. We write $\stp_{k,q}=\stp^{G}_{k,q}$.
\end{definition}

We represent a scatter type canonically as a sorted list of its triples
$(r,\nu,\stp_{k,q}(r,\nu))$, using the canonical representation of $\nu$.
Its representation size, denoted by $\size{\stp_{k,q}}$, is the total size of
this list, including the encodings of the radii and values.
% lax end

% lax begin Lax185574.ScatterTypeRestriction
\begin{lemma}\label{obs:stp-restrict}
For every partial graph $G$ and $q\ge 1$, the function $\stp_{k+1,q-1}$ is the restriction of $\stp_{k,q}$ to $\operatorname{domain}(\stp_{k+1,q-1})$.
\end{lemma}
% lax end
% lax begin Lax185574Proofs.ScatterTypeRestriction.stp_restrict
\begin{proof}
Level $i$ of $\operatorname{domain}(\stp_{k+1,q-1})$ equals level $i+1$ of $\operatorname{domain}(\stp_{k,q})$: the ranks $\bigl((k+1)+i,(q-1)-i\bigr)=\bigl(k+(i+1),q-(i+1)\bigr)$ and the radius bounds coincide, the value caps $(k+1)+(q-1)=k+q$ coincide, and both functions use the same canonical set $S_r([\nu]^G)$.
\end{proof}
% lax end

In the model checking application, the type trees of all vertices are maintained by a dynamic-programming algorithm, and the scatter type is computed from them.
Recall that $|G|=|V(G)|$ and $\size G=|V(G)|+|R(G)|$.

% lax begin Lax185574.ComputeScatterType
\begin{lemma}\label{lem:compute-stp}
There is an algorithm which, given a nonempty partial graph $G$, numbers
$k\ge1$ and $q\in\N$, its distance matrix $D_G$, and the family
$(\tree_{k,q}(a):a\in V(G))$, computes $\stp^G_{k,q}$ in time
\[
|G|^2\cdot
\poly_{k,q}\left(\max_{a\in V(G)}\size{\tree_{k,q}(a)}\right).
\]
\end{lemma}
% lax end
% lax begin Lax185574Proofs.ComputeScatterType.exists_program_stp
\begin{proof}
Put $M\coloneqq\max_{a\in V(G)}\size{\tree_{k,q}(a)}$.
For every $a\in V(G)$ and $i\in[q]$, applying \Cref{lem:shift} $i$ times to $\tree_{k,q}(a)$ yields $\tree_{k+i,q-i}(a)$, of size $\poly_{k,q}(M)$.
This takes $|G|\cdot\poly_{k,q}(M)$ time and gives $\operatorname{domain}(\stp_{k,q})$, as each radius is $O_{k,q}(1)$; in particular, $|\operatorname{domain}(\stp_{k,q})|=O_{k,q}(|G|)$.

For each $(r,\nu)\in \operatorname{domain}(\stp_{k,q})$, run the greedy process defining $S_r([\nu]^G)$, stopping after $k+q$ vertices were discovered.
It scans $V(G)$ in order, testing membership in $[\nu]^G$ by comparing the trees computed above with $\nu$. After adding a vertex $v$, scan the row $D_G(v,\cdot)$ and mark all vertices $u$ with $D_G(v,u)\le r$ as unavailable.
At most $k+q$ rows are scanned, so one run takes $|G|\cdot\poly_{k,q}(M)$ time, and the $|\operatorname{domain}(\stp_{k,q})|=O_{k,q}(|G|)$ runs together take the stated time.
\end{proof}
% lax end

\subsection{The far tree}\label{sec:ftree}

% lax begin Lax185574.FarTrees
\begin{definition}\label{def:ftree}
For a partial graph $G$, $k\in\N$, $q\ge1$, and a nonempty tuple $\tup a\in V(G)^{\le k}$, define
\[
\ftree_{k,q}(\tup a)\coloneqq
\setof{\tree_{k+1,q-1}(b)}{b\in V(G),\ \dist(\tup a,b)>\rhom(k+1,q-1)}.
\]
\end{definition}
% lax end

The goal of this subsection is to compute $\ftree_{k,q}(\tup a)$ from $\tree_{k,q}(\tup a)$ and $\stp_{k,q}$ (\Cref{lem:to-ftree}).
% lax begin Lax185574Proofs.Clustering.clustering
Our main tool is the following clustering lemma~\cite[Corollary~3.4 and Lemma~3.5]{locality-note}.
% lax end

% lax begin Lax185574.Clustering
\begin{lemma}[{\cite{locality-note}}]\label{lem:clustering}
Let $r\in\N$, let $\ell\ge1$, let $G$ be a graph, and let $\tup a=(a_1,\ldots,a_\ell)\in V(G)^\ell$.
There exist a radius $\crad$ with $r\le \crad\le9^{\ell-1}r$ and a partition $\mathcal B$ of $\set{a_1,\ldots,a_\ell}$ with the following property.

For every $X\subseteq V(G)$, inclusion-wise maximal $4\crad$-scattered subset $S \subseteq X$, and $H\ge4\crad+r$,
\[
\exists x\in X\colon \dist^G(x,\tup a)>r
\quad\Longleftrightarrow\quad
\begin{aligned}[t]
&\exists x\in X\colon r<\dist^G(x,\tup a)\le H,\\[-2pt]
&\text{or}\quad \bigl|\setoff{B\in\mathcal B}{\exists x\in X\colon \dist^G(x,B)\le r}\bigr|< |S|.
\end{aligned}
\]
Moreover, the radius $\crad$ and the partition $\mathcal B$ can be determined by tests of the form $\dist^G(a_i,a_j)\le s$, where $i,j\in[\ell]$ and $s\le8\cdot9^{\ell-1}r$.
\end{lemma}
% lax end

% lax begin Lax185574.ToFarTree
\begin{lemma}\label{lem:to-ftree}
Fix $k\in\N$ and $q\ge1$, let $G$ be a partial graph, and let $\tup a\in V(G)^{\le k}$ be nonempty.
Then $\ftree_{k,q}(\tup a)$ is computable from $\tree_{k,q}(\tup a)$ and $\stp_{k,q}$ in time
\[
\size{\stp_{k,q}}\cdot
\poly_{k,q}\bigl(\size{\tree_{k,q}(\tup a)}\bigr).
\]
\end{lemma}
% lax end
% lax begin Lax185574Proofs.ToFarTree.exists_program_ftree
\begin{proof}
Let
$\tau\coloneqq\tree_{k,q}(\tup a)$,
$\sigma\coloneqq\stp_{k,q}$,
$r\coloneqq\rhom(k+1,q-1)$,
$H\coloneqq\rhop(k+1,q-1)$, and
$\ell\coloneqq|\tup a|$.
The entries with $i=1$ in the domain of $\sigma$ let us compute the set
\[
\mathcal T
\coloneqq
\setoff{\tree_{k+1,q-1}(b)}{b\in V(G)}.
\]
Then
\[
\ftree_{k,q}(\tup a) =\setoff{\nu\in\mathcal T}{\exists b\in[\nu]^G\colon\dist(\tup a,b)>r}.
\]

\paragraph{Rewriting the target.}
The tree $\tau$ only records vertices within distance $H$ of $\tup a$, while a vertex $b\in[\nu]^G$ with $\dist(\tup a,b)>r$ may be farther away than $H$.
We will now use \Cref{lem:clustering} to replace the existence of such a vertex by two conditions that we can check.
Apply the lemma with radius $r$ to the Gaifman graph of $G$ (with vertex set $V(G)$ and edge set $R(G)$) and to $\tup a$, and obtain a radius $\crad$ and a partition $\mathcal B$ of $\set{a_1,\ldots,a_\ell}$.
Since $\ell\le k$,
\[
4\crad\le4\cdot9^{\ell-1}r\le9^{k+1}r,
\qquad
4\crad+r\le(4\cdot9^{\ell-1}+1)r\le9^{k}r\le H.
\]
The first chain shows that $(4\crad,\nu)$ lies in the domain of $\sigma$ for every $\nu\in\mathcal T$. The second is the hypothesis $H\ge4\crad+r$ of \Cref{lem:clustering}, so the lemma applies to every $X\subseteq V(G)$ together with $S_{4\crad}(X)$.
Define
\begin{align*}
\mathcal N&\coloneqq\setof{\tree_{k+1,q-1}(b)}{b\in V(G),\ r<\dist(\tup a,b)\le H},\\
\mathcal N_B&\coloneqq\setof{\tree_{k+1,q-1}(b)}{b\in V(G),\ \dist(b,B)\le r}
\qquad\text{for $B\in\mathcal B$.}
\end{align*}
For $\nu\in\mathcal T$, \Cref{lem:clustering} applied to $X\coloneqq[\nu]^G$ and $S\coloneqq S_{4\crad}([\nu]^G)$ states that
\[
\exists b\in[\nu]^G\colon\dist(\tup a,b)>r
\quad\Longleftrightarrow\quad
\nu\in\mathcal N\ \text{ or }\ \bigl|\setoff{B\in\mathcal B}{\nu\in\mathcal N_B}\bigr|<|S_{4\crad}([\nu]^G)|.
\]
As $\sigma(4\crad,\nu)=\min(|S_{4\crad}([\nu]^G)|,k+q)$ by \Cref{def:stp}, and $\bigl|\setoff{B\in\mathcal B}{\nu\in\mathcal N_B}\bigr|\le|\mathcal B|\le\ell\le k<k+q$, we can equivalently replace the bound $|S_{4\crad}([\nu]^G)|$ above by $\sigma(4\crad,\nu)$.
Therefore,
\[
\ftree_{k,q}(\tup a)
=\setoff{\nu\in\mathcal T}{\nu\in\mathcal N\ \text{ or }\ \bigl|\setoff{B\in\mathcal B}{\nu\in\mathcal N_B}\bigr|<\sigma(4\crad,\nu)}.
\]
In this expression, $\mathcal T$ and the values $\sigma(4\crad,\nu)$ can be read from $\sigma$.
It remains to compute $\crad$, $\mathcal B$, $\mathcal N$, and the sets $\mathcal N_B$ from $\tau$.

\paragraph{Reading $\crad$ and $\mathcal B$ off $\tau$.}
By \Cref{lem:clustering}, $\crad$ and $\mathcal B$ are determined by tests $\dist(a_i,a_j)\le s$ with $i,j\in[\ell]$ and
\[
s\le8\cdot9^{\ell-1}r\le9^kr\le H\le\rhom(k,q).
\]
These are binary distance atoms of radius at most $\rhom(k,q)$ and hence determined by the label $\atp_{k,q}(\tup a)$ of the root of $\tau$.
We represent every set $B\in\mathcal B$ by its set of positions $I_B\coloneqq\setof{i\in[\ell]}{a_i\in B}$, which is determined by the equality atoms in the root atomic type of $\tau$.

\paragraph{Reading $\mathcal N$ and $\mathcal N_B$ off $\tau$.}
The trees attached to the root of $\tau$ are exactly
the trees $\tree_{k+1,q-1}(\tup ab)$ with $b\in V(G)$ and $\dist(\tup a,b)\le H$.
Such a tree $\mu$ determines two things about $b$: by \Cref{lem:restrict}, $\mu\restriction{(\ell+1)}=\tree_{k+1,q-1}(b)$, and since the root atomic type $\atp_{k+1,q-1}(\tup ab)$ of $\mu$ contains the binary distance atoms of radius at most $\rhom(k+1,q-1)=r$, we have $\mu\models\dist(x_{j},x_{\ell+1})\le r$ if and only if $\dist(a_{j},b)\le r$, for every $j\in[\ell]$.
Therefore,
\[
\mathcal N
=\setof{\mu\restriction{(\ell+1)}}{(i,\mu)\in\mathcal E_{k,q}(\tup a),\ \mu\models\textstyle\bigwedge_{j\in[\ell]}\dist(x_{j},x_{\ell+1})>r}.
\]
Likewise, since $r\le H$, the condition $\dist(b,B)\le r$ implies $\dist(\tup a,b)\le H$, and
\[
\mathcal N_B
=\setof{\mu\restriction{(\ell+1)}}{(i,\mu)\in\mathcal E_{k,q}(\tup a),\ \mu\models\textstyle\bigvee_{j\in I_B}\dist(x_{j},x_{\ell+1})\le r}.
\]
The right-hand sides depend only on $\tau$.
Substituting them into the expression for $\ftree_{k,q}(\tup a)$ above gives a description in terms of $\tau$ and $\sigma$ only.

\paragraph{Running time.}
The radius $\crad$ and the sets $I_B$ are read off the root label of $\tau$.
For all extensions $(i,\mu)$ of $\tau$, the trees $\mu\restriction({\ell+1})$ are obtained by \Cref{lem:restrict} and the atoms $\dist(x_{j},x_{\ell+1})\le r$ are read off the root of $\mu$, in time $\poly_{k,q}(\size\tau)$ in total.
Collecting them yields $\mathcal N$ and the at most $k$ sets $\mathcal N_B$ within the same bound.
The expression for $\ftree_{k,q}(\tup a)$ is then evaluated by one scan of $\sigma$: for every entry $(4\crad,\nu,\sigma(4\crad,\nu))$ with $\nu$ of rank $(k+1,q-1)$, membership of $\nu$ in $\mathcal N$ and in each $\mathcal N_B$ is tested by comparing $\nu$ with every element of these at most $k+1$ sets, in time $O(k\cdot\size\tau\cdot(\size\nu+\size\tau))$.
Summed over all entries of $\sigma$ this is $O(\size\sigma\cdot\poly_{k,q}(\size\tau))$, which gives the stated running time.
\end{proof}
% lax end

\section{Lifting type trees through merges and resolves}
\label{sec:merge-resolve}
The model checking algorithm will process
the construction sequence step by step and maintain the type trees of all vertices.
We first show how to lift type trees through a merge, and then through a resolve.

% lax begin Lax185574.MergeLift
\begin{lemma}[merge lift]\label{lem:merge-lift}
Fix $\ell\le k$ and $q\in\N$.
Let $G$ be a partial graph and $\tup a\in V(G)^\ell$.
Given
$\tree^G_{k,q}(\tup a)$ and two distinct colors $P,Q$, one can compute
 $\tree^{H}_{k,q}(\tup a)$, where $H=\merge PQ(G)$,
in time
\[
\poly_{k,q}(\size{\tree^G_{k,q}(\tup a)}).
\]
\end{lemma}
% lax end
% lax begin Lax185574Proofs.MergeLift.exists_program_mergeLift
\begin{proof}
We proceed by induction on $q$.

\paragraph{Root labels.}
The partial graphs $G$ and $H$ have the same resolved pairs, so all distances agree, and only the colors $P$ and $Q$ change.
Hence, for every vertex $u$ and $s\in\N$,
\begin{align*}
H\models P(u)&\quad\Longleftrightarrow\quad G\models P(u)\lor Q(u),\\
H\models Q(u)&\quad\Longleftrightarrow\quad\textit{false},\\
H\models\dist(u,P)<s&\quad\Longleftrightarrow\quad G\models(\dist(u,P)<s)\lor(\dist(u,Q)<s),\\
H\models\dist(u,Q)<s&\quad\Longleftrightarrow\quad\textit{false},
\end{align*}
and every other atom has the same interpretation in $G$ and $H$.
This allows us to rewrite every atom \(\phi\)
into a quantifier-free formula \(\phi'\)
of the same distance rank, so that
$H\models\phi(\tup a)$ if and only if $G\models\phi'(\tup a)$.
Since whether or not $G\models\phi'(\tup a)$ is determined by $\atp^G_{k,q}(\tup a)$, it follows that the new root label $\atp^{H}_{k,q}(\tup a)$ can be computed  from $\atp^G_{k,q}(\tup a)$.
Only atoms mentioning $P$ or $Q$ change, so the new root label is obtained by one pass over the old root label and is no larger.
This settles the root label.

\paragraph{Extensions.}
Let $q\ge1$.
Since $\dist^{H}=\dist^G$, every guard admits the same vertices in $G$ and $H$, so the inductive assumption gives $\mathcal E^{H}_{k,q}(\tup a)=\setof{(i,F(\mu))}{(i,\mu)\in\mathcal E^G_{k,q}(\tup a)}$,
where $F(\mu)$ is obtained inductively.
Together with the root label, this gives $\tree^{H}_{k,q}(\tup a)$.

\paragraph{Running time.}
The transformation replaces atoms in the label of every node and keeps the guards, so it visits each node of the input tree once, and no tree grows.
By \Cref{obs:canonical}, restoring canonical form takes time polynomial  in the size of the tree.
\end{proof}
% lax end

A resolve can bring vertices far from a tuple within guard range.
Their type trees are not recorded in the old type tree of the tuple, and we obtain them from the far tree instead.

% lax begin Lax185574.ResolveLift
\begin{lemma}[resolve lift]\label{lem:resolve-lift}
Let $k\ge1$ and $q\in\N$.
Given a partial graph $G$, its distance matrix $D_G$, colors $P,Q$,
a sign $\pm$, and
$(\tree^G_{k,q}(a):a\in V(G))$, one can compute
$H=\resolve\pm PQ(G)$, its distance matrix $D_H$, and
$(\tree^{H}_{k,q}(a):a\in V(G))$ in time
\[
|G|^2\cdot\poly_{k,q}(M+T),
\]
where
\begin{align*}
M&\coloneqq\max_{a\in V(G)}\size{\tree^G_{k,q}(a)},\\
T&\coloneqq\max_{a\in V(G)}\bigl|\bigl\{\tree^G_{k+i,q-i}(\tup v)\colon 0\le i\le q,\ \ell\in[k+i],\ \tup v\in\bigl(\Ball^{H}_{\rhop(k,q)}(a)\bigr)^{\ell}\bigr\}\bigr|.
\end{align*}
Each output tree has size $\poly_{k,q}(T)$.
\end{lemma}
% lax end
% lax begin Lax185574Proofs.ResolveLift.exists_program_resolveLift
\begin{proof}
Assume the resolve is positive, as the negative case is symmetric.
By induction on $q$,
we construct an auxiliary map $F$ on trees of arbitrary nonempty tuples, so that $\tree^{H}_{k,q}(\tup a)=F(\tree^G_{k,q}(\tup a))$ for all $\ell\le k$ and $\tup a\in V(G)^\ell$.
Fix $\tup a \in V(G)^\ell$ and
put $\tau\coloneqq\tree^G_{k,q}(\tup a)$ and $\sigma\coloneqq\stp^G_{k,q}$.

\paragraph{Root labels.}
We may assume that $P^G$ and $Q^G$ are nonempty, as otherwise $H=G$.
We rewrite every atom $\phi$ in $H$ to an equivalent formula $\phi'$ in $G$,
so that $\phi'$ has the same distance rank as $\phi$,
and $H\models\phi(\tup a)\iff G\models\phi'(\tup a)$.
Below, $u,v$ are vertices, $Y$ is a color, and $s\in\N$.

\emph{First-order atoms.}
Equality, color, and non-edge atoms $N(x,y)$ have the same interpretation in $G$ and~$H$.
By definition of the positive resolve,
\[
H\models E(u,v)
\quad\Longleftrightarrow\quad
G\models E(u,v)\lor\Bigl(\bigl((P(u)\land Q(v))\lor(Q(u)\land P(v))\bigr)\land u\ne v\land\lnot N(u,v)\Bigr).
\]

\emph{Distances.}
We have $R(H)=R(G)\cup P^GQ^G$.
A shortest path from $u$ to $v$ in $H$ is a path in $G$, or runs in $G$ from $u$ to $P^G$ or $Q^G$, then crosses to the other part or returns to the same part through the other one, and then runs in $G$ to $v$:
otherwise, the segment between its first and last vertex in $P^G\cup Q^G$ could be replaced by one pair from $P^GQ^G$, or by two through the other part if both ends lie in the same part.
Hence
\begin{equation}\label{eq:resolve-dist}
\dist^H(u,v)=\min\left\{
\begin{aligned}
&\dist^G(u,v),\\
&\dist^G(u,P)+1+\dist^G(Q,v),\quad \dist^G(u,Q)+1+\dist^G(P,v),\\
&\dist^G(u,P)+2+\dist^G(P,v),\quad \dist^G(u,Q)+2+\dist^G(Q,v)
\end{aligned}\right\}.
\end{equation}

\emph{Unary distance atoms.}
Taking the minimum of \eqref{eq:resolve-dist} over $v\in Y^G$ gives
\begin{align*}
&H\models\dist(u,Y)<s
\quad\Longleftrightarrow\quad
G\models\dist(u,Y)<s\\
&\qquad\lor\dist(u,P)<(s-1-\dist^G(Q,Y))\lor\dist(u,Q)<(s-1-\dist^G(P,Y))\\
&\qquad\lor\dist(u,P)<(s-2-\dist^G(P,Y))\lor\dist(u,Q)<(s-2-\dist^G(Q,Y)).
\end{align*}
The values $\dist^G(P,Y)$ and $\dist^G(Q,Y)$ do not depend on $u$, and are computed once. A unary-distance comparison whose upper bound is nonpositive, or from which an infinite distance is subtracted, is interpreted as \emph{false}.

\emph{Binary distance atoms.}
Define
\begin{align*}
&\delta^s_{P,Q}(x,y)\coloneqq\\
&\bigvee_{s_1+1+s_2=s}\bigl(\dist(x,P)<s_1{+}1\land\dist(y,Q)<s_2{+}1\bigr)\lor\bigl(\dist(x,Q)<s_1{+}1\land\dist(y,P)<s_2{+}1\bigr)\\
\lor{}&\bigvee_{s_1+2+s_2=s}\bigl(\dist(x,P)<s_1{+}1\land\dist(y,P)<s_2{+}1\bigr)\lor\bigl(\dist(x,Q)<s_1{+}1\land\dist(y,Q)<s_2{+}1\bigr).
\end{align*}
By \eqref{eq:resolve-dist},
\begin{align}\label{eq:bin-dist}
H\models\dist(u,v)\le s
\quad\Longleftrightarrow\quad
G\models(\dist(u,v)\le s)\lor\delta^s_{P,Q}(u,v).
\end{align}

\emph{Root label.}
For an atom \(\phi\), let $\phi'$ be the corresponding formula as obtained above.
Then $H\models\phi(\tup a)$ if and only if $G\models\phi'(\tup a)$ and $\phi'$ has the same distance rank as $\phi$.
Since whether or not $G\models\phi'(\tup a)$ is determined by $\atp^G_{k,q}(\tup a)$, it follows that the new root label $\atp^{H}_{k,q}(\tup a)$ can be computed  from $\atp^G_{k,q}(\tup a)$.
Algorithmically, there are $O_{k,q}(1)$ atoms without colors, and each nonempty color contributes $O_{k,q}(1)$ further atoms.

\paragraph{Extensions.}
Let $q\ge1$, and put $r\coloneqq\rhom(k+1,q-1)$ and $R\coloneqq\rhop(k+1,q-1)$.
On trees of rank $(k+1,q-1)$, $F$ denotes the map given by the inductive assumption, using the restriction of $\sigma$ to $\operatorname{domain}(\stp_{k+1,q-1})$ from \Cref{obs:stp-restrict} and the same precomputed distances.
By the inductive assumption,
\begin{align}\label{neqE}
\mathcal E^{H}_{k,q}(\tup a)=\setof{\bigl(i,F(\tree^G_{k+1,q-1}(\tup a b))\bigr)}{i\in[\ell], \dist^H(a_i,b)\le R}.
\end{align}

For $i\in[\ell]$, put
\[
\gamma_i(x_1,\ldots,x_{\ell+1})\coloneqq\delta^R_{P,Q}(x_i,x_{\ell+1}).
\]
By \eqref{eq:bin-dist},
\begin{align}\label{eq:Disteq}
H\models\dist(a_i,b)\le R
\quad\Longleftrightarrow\quad
G\models\dist(a_i,b)\le R\ \lor\ \gamma_i(\tup a b).
   \end{align}
The right-hand side of \eqref{eq:Disteq}
is equivalent to the disjunction of three conditions:
\begin{enumerate}
   \item  $\dist^G(a_i,b)\le R$,
   \item  $\dist^G(\tup a,b)\le R$ and $G\models\gamma_i(\tup a b)$,
    \item  $\dist^G(\tup a,b)>r $ and $G\models\gamma_i(\tup a b)$.
\end{enumerate}
This holds since $r\le R$ by \eqref{eq:horizon}.
Consequently,
\[
\mathcal E^{H}_{k,q}(\tup a)=\mathcal E_{\mathrm{old}}\cup\mathcal E_{\mathrm{near}}\cup\mathcal E_{\mathrm{far}},
\]
where
\begin{align*}
\mathcal E_{\mathrm{old}}&\coloneqq\setof{\bigl(i,\ F(\tree^{G}_{k+1,q-1}(\tup a b))\bigr)}{i\in[\ell],\ b\in V(G),\ \dist^G(a_i,b)\le R},\\
\mathcal E_{\mathrm{near}}&\coloneqq\setof{\bigl(i,\ F(\tree^{G}_{k+1,q-1}(\tup a b))\bigr)}{i\in[\ell],\ b\in V(G),\ \dist^G(\tup a,b)\le R,\ G\models\gamma_i(\tup a b)},\\
\mathcal E_{\mathrm{far}}&\coloneqq\setof{\bigl(i,\ F(\tree^{G}_{k+1,q-1}(\tup a b))\bigr)}{i\in[\ell],\ b\in V(G),\ \dist^G(\tup a,b)>r,\ G\models\gamma_i(\tup a b)}.
\end{align*}

The first set is obtained via
\[
\mathcal E_{\mathrm{old}}
=
\setof{(i,F(\mu))}{(i,\mu)\in\mathcal E^G_{k,q}(\tup a)}.
\]
Every witness $b$ for $\mathcal E_{\mathrm{near}}$ is admitted by some guard $j$ in $G$, so $\tree^G_{k+1,q-1}(\tup a b)$ occurs as an extension $(j,\mu)$ of $\tau$ with $\mu\models\gamma_i$.
Conversely, every such $\mu$ equals $\tree^G_{k+1,q-1}(\tup a b)$ for some $b$ with $\dist^G(a_j,b)\le R$ and $G\models\gamma_i(\tup a b)$.
Hence, using the notation $\mu\models\gamma$ introduced after \Cref{lem:eval},
\[
\mathcal E_{\mathrm{near}}
=
\setof{(i,F(\mu))}{i\in[\ell],\ (j,\mu)\in\mathcal E^G_{k,q}(\tup a),\ \mu\models\gamma_i}.
\]
For $\mathcal E_{\mathrm{far}}$, consider $b$ with $\dist^G(\tup a,b)>r$.
Since $r=\rhom(k+1,q-1)$ is exactly the separation required by \Cref{lem:compose} at rank $(k+1,q-1)$, that lemma
gives
\[
\tree^G_{k+1,q-1}(\tup a b)
=
\tree^G_{k+1,q-1}(\tup a)\ast\tree^G_{k+1,q-1}(b),
\]
where $\tree^G_{k+1,q-1}(\tup a)$ is obtained from $\tau=\tree^G_{k,q}(\tup a)$ by \Cref{lem:shift}.
As $b$ ranges over all such vertices, $\tree^G_{k+1,q-1}(b)$ ranges exactly over $\ftree^G_{k,q}(\tup a)$, computed from $\tau$ and $\sigma$ by \Cref{lem:to-ftree}.
Therefore
\[
\mathcal E_{\mathrm{far}}
=
\setof{\bigl(i,\ F(\tree^G_{k+1,q-1}(\tup a)\ast\nu)\bigr)}{i\in[\ell],\ \nu\in\ftree^G_{k,q}(\tup a),\ \tree^G_{k+1,q-1}(\tup a)\ast\nu\models\gamma_i}.
\]
So the new root label, $\mathcal E_{\mathrm{old}}$, $\mathcal E_{\mathrm{near}}$, and $\mathcal E_{\mathrm{far}}$ are computed from $\tau$ together with $\sigma$ and the values $\dist^G(P,Y)$ and $\dist^G(Q,Y)$, which do not depend on $\tup a$. This defines $F(\tau)=\tree^{H}_{k,q}(\tup a)$.

\paragraph{Running time.}
Computing $H$ takes time $O(\size{H})=O(|G|^2)$.
Scanning $D_G$ gives $\dist^G(u,P)$ and $\dist^G(u,Q)$ for all vertices $u$, and thus $\dist^G(P,Y)$ and $\dist^G(Q,Y)$ for all nonempty colors $Y$, in time $O(|G|^2)$.
Then \eqref{eq:resolve-dist} gives $D_H$ in time $O(|G|^2)$, and \Cref{lem:compute-stp} computes $\sigma\coloneqq\stp^G_{k,q}$ from $D_G$ in time $|G|^2\cdot\poly_{k,q}(M)$.
Every tree $\tree^G_{k+i,q-i}(b)$ is obtained from $\tree^G_{k,q}(b)$ by at most $q$ applications of \Cref{lem:shift}, so it has size $\poly_{k,q}(M)$, and $\size\sigma\le |G|\cdot\poly_{k,q}(M)$, since the domain of $\sigma$ has $O_{k,q}(1)$ entries per vertex.

\emph{Number of calls.}
Fix $a\in V(G)$ and consider the recursion computing $F(\tree^G_{k,q}(a))$.
We collect the distinct old extension trees admitted by at least one guard in $H$ before recursing, retaining the list of guards for each tree.
The calls at depth $j$ are on trees $\tree^G_{k+j,q-j}(\tup v)$ for tuples $\tup v$ reached by the guards in $H$.
By \Cref{lem:locality} applied to $H$, every such $\tup v$ is a tuple from $\Ball^{H}_{\rhop(k,q)}(a)$.
So by definition of $T$, a call makes at most $T$ recursive calls, and as the depth is $q$, there are $\poly_{k,q}(T)$ calls for $a$.

\emph{Cost per call and output size.}
A recursive input is a child of the previous input, or is obtained by rank shift and far composition with a  tree of arity one recorded in $\sigma$.
These trees have size $\poly_{k,q}(M)$, so by \Cref{lem:shift,lem:compose} and the depth bound $q$, every recursive input and every candidate extension tree has size $\poly_{k,q}(M)$.
For a call with positive remaining depth, restricting $\sigma$ and applying \Cref{lem:to-ftree} takes time $O(\size\sigma\cdot\poly_{k,q}(M))=O(|G|\cdot\poly_{k,q}(M))$.
The far tree has at most $|G|$ members, so composing them and testing the formulas $\gamma_i$
(which are quantifier-free formulas of distance rank $(k+1,q-1)$, so are determined by the root label of a tree of rank $(k+1,q-1)$) takes the same time.
We scan the candidates, test their guards, and compare each admitted candidate with the list of at most $T$ distinct admitted trees collected so far.
Including the candidates from the old extensions, this takes $O(|G|\cdot\poly_{k,q}(M+T))$ time.

The colors recorded in each new root label meet $\Ball^{H}_{\rhop(k,q)}(a)$ by \Cref{lem:locality} applied to $H$, and as vertices of distinct colors have distinct old trees, there are at most $T$ of them.
So each new root label has size $O_{k,q}(T)$ and is computed in time $O(|G|\cdot\size\tau)$.
Together with the branching and depth bounds, this gives output size $\poly_{k,q}(T)$; restoring canonical form has the same polynomial bound by \Cref{obs:canonical}.

Altogether the recursion for $a$ takes time $O(|G|\cdot\poly_{k,q}(M+T))$, and for all vertices $O(|G|^2\cdot\poly_{k,q}(M+T))$.
Adding the construction of $H$ and $D_H$, the distance preprocessing, and the scatter type gives the claimed bound.
\end{proof}
% lax end

\section{The model checking algorithm}\label{sec:model-checking}
We now give the model checking algorithm,
proving \Cref{thm:mc}.
Recall the width notion used in the statement.
\defrqwidth*

The width bounds numbers of ordinary types, and the following lemma converts such bounds into bounds on type trees, as required by \Cref{lem:resolve-lift}.

% lax begin Lax185574.TypeTreeSizes
\begin{lemma}[type tree sizes]\label{lem:ball-tree-bound}
Let $k\ge1$ and $q\in\N$, and put $r\coloneqq\rhop(k,q)$ and $d\coloneqq r+k+q$.
The following holds for every partial graph $G$.
\begin{enumerate}[label=(\roman*)]
\item\label{lem:ball-count}
$\bigl|\setoff{\tree^G_{k+i,q-i}(\tup v)}{\tup v\in X^{\ell}}\bigr|
\le|\Types^G_{d,d}(X)|$ for all $i\in\set{0,\ldots,q}$, $\ell\in [k+i]$, and $X\subseteq V(G)$.
\item\label{lem:ball-size}
$\size{\tree^G_{k,q}(a)}
\le
\poly_{k,q}\bigl(|\Types^G_{d,d}(\Ball^G_r(a))|\bigr)$ for every $a\in V(G)$.
\end{enumerate}
\end{lemma}
% lax end
\begin{proof}
We first show that ordinary types determine type trees: for all $k',q'\in\N$, every partial graph $G$, and all nonempty tuples $\tup a,\tup b\in V(G)^{\le k'}$ of the same length,
\begin{equation}\label{eq:tp-to-tree}
\tp^G_{\rhop(k',q')+q'}(\tup a)=\tp^G_{\rhop(k',q')+q'}(\tup b)
\quad\Longrightarrow\quad
\tree^G_{k',q'}(\tup a)=\tree^G_{k',q'}(\tup b).
\end{equation}
By expressing distance atoms by paths along $E\lor N$, every \distFO formula of distance rank $(k',q')$ is equivalent to a first-order formula of quantifier rank at most $\rhop(k',q')+q'$. Thus, the left-hand side of \eqref{eq:tp-to-tree} implies that $\tup a$ and $\tup b$ satisfy
the same \distFO formulas of distance rank $(k',q')$.
By \Cref{cor:tp-tree}, this yields the right-hand side.

% lax begin Lax185574Proofs.TypeTreeSizes.treeCount_le
For \ref{lem:ball-count}, extend every tuple $\tup v\in X^\ell$ to a tuple $\tup v'\in X^d$ by repeating its last coordinate.
Then $\tp_d(\tup v')$ determines $\tp_d(\tup v)$, which determines $\tree_{k+i,q-i}(\tup v)$ by \eqref{eq:tp-to-tree}, as $\rhop(k+i,q-i)+q-i\le \rhop(k,q)+q\le d$ by \eqref{eq:horizon-cons}.
So distinct type trees yield distinct elements of $\Types^G_{d,d}(X)$.
% lax end

% lax begin Lax185574Proofs.TypeTreeSizes.tree_length_le
For \ref{lem:ball-size}, put $X\coloneqq\Ball_r(a)$ and $K\coloneqq|\Types^G_{d,d}(X)|$.
As every vertex has exactly one color, vertices of distinct colors have distinct type trees, so at most $K$ colors meet $X$ by \ref{lem:ball-count}.
By \Cref{lem:locality}, every node of $\tree_{k,q}(a)$ is $\tree_{k+j,q-j}(\tup v)$ for a tuple $\tup v$ from $X$, and every color in a true atom of its root meets $X$.
So every root has size $O_{k,q}(K)$, and by \ref{lem:ball-count}, every node has at most $K$ children per guard, and thus at most $(k+q)K$ children.
As the depth is $q$, this gives $\size{\tree_{k,q}(a)}\le\poly_{k,q}(K)$.
% lax end
\end{proof}

{\renewcommand{\mcdepth}{d}\mcthm*}
% lax begin Lax185574Proofs.ConstructionModelChecking.exists_program_modelChecking_sequence
\begin{proof}
Let $q$ be the quantifier rank of $\phi$, and let $r,d$ be as in
\Cref{lem:ball-tree-bound} for $k=1$ and this $q$.  We prove the statement
with radius $r$ and depth $d$.
If $n\le1$, evaluate $\phi$ directly, so assume $n\ge2$.

Let $G_1,\ldots,G_\lambda$ be the given
construction sequence of $G$, and let $K$ be its radius-$r$ depth-$d$
width.  We iteratively compute $(\tree^{G_t}_{1,q}(a):a\in V(G))$ for $t=1,\ldots,\lambda$. To save time we also maintain the distance matrix $D_{G_t}$.
Initially, for \(t=1\), there are no resolved pairs and each nonempty part is a singleton.
Hence, every guard admits only vertices of the guarded tuple itself, so
each \(\tree^{G_1}_{1,q}(a)\) is determined by the part containing \(a\), and can therefore be computed in time \(O_q(1)\).
Initialize also $D_{G_1}$ in $O(n^2)$ time by putting $0$ on its diagonal and $\infty$ elsewhere.

Now let \(t<\lambda\), and assume the trees in $G_t$ have been computed.
All these trees have size $\poly_{q}(K)$: for every $t'\in[\lambda-1]$
and $a\in V(G)$, \Cref{lem:ball-tree-bound}\ref{lem:ball-size} bounds
$\size{\tree^{G_{t'}}_{1,q}(a)}$ by a polynomial in
$|\Types^{G_{t'}}_{d,d}(\Ball^{G_{t'}}_{r}(a))|$.  Using $R(G_{t'})\subseteq R(G_{t'+1})$
and therefore $\Ball^{G_{t'}}_r(a)\subseteq \Ball^{G_{t'+1}}_r(a)$, this is at most
$|\Types^{G_{t'}}_{d,d}(\Ball^{G_{t'+1}}_{r}(a))|\le K$ by the definition of
the width.

For a merge, construct $G_{t+1}$ in time $O(n)$, keep $D_{G_{t+1}}=D_{G_t}$ unchanged, and apply \Cref{lem:merge-lift} to each tree of arity one, taking $O(n\cdot\poly_q(K))$ time in total.
For a resolve, apply \Cref{lem:resolve-lift} to $G_t$, the supplied matrix $D_{G_t}$, and all its trees of arity one. The lemma returns both $D_{G_{t+1}}$ and the new trees.
Its parameter $M$ is $\poly_q(K)$ by the bound above.
Its parameter $T$ is
$$T\coloneqq \max_{a\in V(G)}\left|\setof{\tree_{1+i,q-i}(\tup v)}{\tup v\in \left(\Ball^{G_{t+1}}_{\rho^+(1,q)}(a)\right)^\ell,\,  0\le i\le q,\, \ell\in [i+1]}\right|.$$
The tuples $\tup v$ above lie in $X\coloneqq\Ball^{G_{t+1}}_r(a)$ since $\rhop(1,q)\le r$.
By \Cref{lem:ball-tree-bound}\ref{lem:ball-count}, there are at most $|\Types^{G_t}_{d,d}(X)|\le K$ trees for each $i$ and arity $\ell$.
There are at most $(q+1)^2$ choices of $i$ and $\ell$, so $T\le(q+1)^2K$.

As each step takes $O(n^2\cdot\poly_q(K))$ time, the total time is $O(\lambda n^2\cdot\poly_\phi(K))$.
After $\lambda-1$ steps we have all the type trees in $G_\lambda$.
These final trees also have size $\poly_q(K)$: if the last step is a merge, no tree grows; if it is a resolve, use the output-size bound $\poly_q(T)$ in \Cref{lem:resolve-lift} and $T\le(q+1)^2K$.

As all pairs of $G_\lambda$ are resolved,
we can turn each quantifier of $\phi$ into a local one without changing its truth value.
Construct \(\phi^{\mathrm{loc}}(x)\) from \(\phi\) by
introducing a fresh variable $x$, and replacing every subformula
$\exists y\,\psi$ of $\phi$ that lies in the scope of quantifiers
$\exists z_1,\ldots,\exists z_j$ by
\[
\exists y~\Bigl(
\dist(x,y)\le\rhop(j+2,q-j-1)
\land\psi
\Bigr),
\]
where $\psi$ is transformed in the same way.  This yields a \distFO formula
$\phi^{\mathrm{loc}}(x)$ such that, for
every $a\in V(G)$,
\[
G_\lambda\models\phi
\quad\Longleftrightarrow\quad
G_\lambda\models\phi^{\mathrm{loc}}(a).
\]
Note that $\phi^{\mathrm{loc}}$ has distance rank $(1,q)$.

Now choose any $a\in V(G)$,
and use \Cref{lem:eval} for $k=1$ to compute whether $G_\lambda\models\phi^{\mathrm{loc}}(a)$,
and hence whether $G_\lambda\models\phi$, from
$\tree^{G_\lambda}_{1,q}(a)$ in time
$\poly\bigl(|\phi^{\mathrm{loc}}|+\size{\tree^{G_\lambda}_{1,q}(a)}\bigr)\le\poly_\phi(K)$.
As $E(G_\lambda)=E$ and $\phi$ is a sentence over graphs, $G_\lambda\models\phi$ if and only if $G\models\phi$.
This fits the preceding time bound and completes the proof.
\end{proof}
% lax end
